\documentclass[options]{philosophersimprint}
%\usepackage{sectsty}
%\linespread{1.2}
\usepackage{xunicode}
\usepackage{xltxtra}
\defaultfontfeatures{Mapping=tex-text}
\setmainfont[
Ligatures={Common}, Numbers={OldStyle}, Variant=01,
BoldFont=LinLibertine_RB.otf,
ItalicFont=LinLibertine_RI.otf,
BoldItalicFont=LinLibertine_RBI.otf
]{LinLibertine_R.otf}
\setmonofont[Scale=0.8]{DejaVuSansMono.ttf}
\usepackage{stmaryrd}
\usepackage[toc,page]{appendix}
\usepackage{amsmath,amsthm,amssymb}
\usepackage[colorlinks=true, citecolor=blue, linkcolor=blue, hyperfootnotes=false]{hyperref}
\usepackage{nameref}
\usepackage{soul}
\usepackage{frame}
\usepackage[style=apa, natbib=true]{biblatex}
\usepackage{scrextend}
\usepackage{mathrsfs}
%\addtokomafont{labelinglabel}{\sffamily}
\makeatletter
\newcommand\labelitem[2]{%
  \item[#1]\def\@currentlabel{#1}\label{#2}}
\makeatother
\usepackage[inline]{enumitem}
\usepackage{framed, csquotes}
\usepackage{tikz}
\usepackage{scalerel}
\usepackage{titlesec}
\newcommand\halo{{{\circ}\mkern-1mu}}
\usepackage{linguex}
\usepackage{parskip}
\usepackage{fancyhdr}
\pagestyle{plain}
% Diagrams
%\usepackage{pgfplots}
\usepackage{tikz}
\usepackage{qtree}

\usepackage{fitch}



\usepackage{FOL_Yale/notation}

\begin{document}

PHIL 1115: First-order Logic \hfill Yale, Fall 2026\\
Lecture 9, Handout \hfill Dr. Daniel Grimmer

\textbf{1. Natural Deduction: Why?} So far, we've been working with a truth-condition-based semantic picture: the meanings of $\ConjTight$, $\Disj$, $\Cond$, and the rest are given by their truth conditions. Consider:
\begin{itemize}
    \item \textit{Truth Conditional Semantics:} To know the meaning of $\ConjTight$ is to know the precise conditions under which claims of the form $A \Conj B$ are true. (Namely, this happens when $A$ is true and $B$ is true.)
\end{itemize}
But here's a fact worth pausing on: you knew what \textit{``and''} meant long before you walked into this classroom. And nobody showed you a truth table in kindergarten. If you separately knew that $p=$``today is Wednesday'' and that $q=$``logic class is fun'', you'd have no trouble concluding $p \Conj q$, that ``today is Wednesday and logic class is fun''. Moreover, if I told you $r \Conj s =$``it is raining and tomorrow is Thursday'' you'd have no trouble concluding $r$, that ``it is raining''. You already knew how to \textit{use} ``and'' correctly in arguments. You knew this long before anyone gave you a truth condition for it.

This lecture will introduce \textbf{Natural Deduction} which begins with \textit{rules of inference} rather than truth conditions. These are patterns of good reasoning like ``from $A \Conj B$, you may infer $A$''. This approach suggests a different picture of what the logical connectives \textit{mean}.
\begin{itemize}
    \item \textit{Proof-Theoretic Semantics:} To know the meaning of $\ConjTight$ is to know two things: firstly, that from any two claims $A$ and $B$ one can derive the claim $A \Conj B$; and secondly, that from any claim of the form $A \Conj B$ one can derive $A$ (and, moreover, one can also derive $B$).

    These are $\ConjTight$'s \textbf{introduction and elimination rules.} We denote them as $A,B \ProvesND A \Conj B$ \quad and\quad $A \Conj B\ProvesND A$ \quad and\quad $A \Conj B\ProvesND B$.
\end{itemize}
Notice that this definition does not rely on the concept of ``truth'', one only needs to understand how $\ConjTight$ functions in a certain language game. The idea that meaning is constituted by use, rather than read off some independently existing truth conditions, has a name in philosophy: \textbf{inferentialism}. (It's a cousin of the later Wittgenstein's slogan that meaning is use, and it's close to what Gerhard Gentzen had in mind in 1934--35, when he (and independently Stanis\l aw Ja\'skowski) first developed natural deduction: Gentzen held that the introduction rule for a connective is, in effect, its \textit{definition}, and the elimination rule is just what falls out of that definition. The scope-line notation we will actually use descends from Ja\'skowski, later streamlined by Frederic Fitch.)

This lecture will introduce a \textit{new method} for dealing with our language of propositional logic, and, like trees, it works by manipulating symbols rather than by surveying models. But do not conclude that it has turned its back on meaning. It marks a shift in spirit from the trees of Lecture 6: there, each tree rule was \textit{answerable to the semantics}, justified by and designed to track the truth-table meaning of the connective it resolved. Natural deduction reverses that order of explanation. Its rules are not read off any prior truth conditions; the licensed inferences are \textit{themselves} offered as what fixes each connective's meaning. That is a rival account of meaning, not an abstention from one: a \textit{proof-theoretic} semantics rather than a truth-conditional one. As such, it gets its own turnstile, $X \ProvesND A$, which is distinct from that of our tree method, $X \Proves A$, and our table method, $X \Entails A$. The meaning of our new symbol, $X \ProvesND A$, is that $A$ can be derived from $X$ using our natural deduction rules. (We'll pin this down precisely below.)

Exactly as before, this new syntactic approach raises questions of soundness and completeness.  Does $X \ProvesND A$ line up with $X \Entails A$? Does it line up with $X \Proves A$? We will not address this question in this course, but here is the answer: $X \ProvesND A$ if and only if $X \Entails A$. Moreover, $X \ProvesND A$ if and only if $X \Proves A$. In other words, natural deduction is both sound and complete! All three of our methods yield exactly the same answers in every case.

\textbf{2. Natural Derivations using $\ConjTight$.} The basic idea underpinning natural deduction is to codify ``natural inferences'' like ``from $A \Conj B$, you may infer $A$'' using syntactic rules. The goal is to derive formulas from a (sometimes empty) set of premises, using the allowed rules.

Consider an argument with premises $P_1 = p \Conj (q \Disj r)$ and $P_2 = (s \Cond t)\Conj p$ and with a conclusion of $C=(q \Disj r) \Conj (s \Cond t)$. How would this argument go in English? From the first premise, $p \Conj (q \Disj r)$, we can immediately infer that $q \Disj r$. Moreover, from the second premise, $(s \Cond t)\Conj p$, we can immediately infer that $s \Cond t$. But putting these two conclusions together we have that $(q \Disj r) \Conj (s \Cond t)$ as desired.

Here is how this argument would be written up as a \textbf{natural derivation}:
\begin{center}
\begin{minipage}[t]{0.19\textwidth}
The formulas above the horizontal line (numbered 1 and 2) are the premises. The $q \Disj r$ formula on the third line indicates that we have inferred this claim (somehow) from the preceding claims. How? We used the \textbf{$\ConjTight$-Elimination Rule} justified by line 1, denoted $\ConjE, 1$.
\end{minipage}
\hfill
\begin{minipage}[t]{0.25\textwidth}
\vspace{-0.5cm}
\begin{align*}
\begin{nd}
\hypo{1}{p \Conj (q \Disj r)}
\hypo{2}{(s \Cond t)\Conj p}
\have{3}{q \Disj r \qquad\qquad\quad \ \ \ConjE, 1}
\have{4}{s \Cond t \qquad\qquad\quad \ \ \ConjE, 2}
\have{5}{(q \Disj r) \Conj (s \Cond t) \quad \ConjI, 3, 4}
\end{nd}
\end{align*}
\end{minipage}
\end{center}
Here is how it works in general:

\textbf{$\ConjTight$-Elimination Rule:}
\vspace{-0.5cm}
\begin{align*}
    \begin{nd}
        \have[n]{pq}{A \Conj B}
        \have[...]{}{\dots}
        \have[m]{}{A\qquad\qquad \ConjE,n}
    \end{nd}
    \hspace{2cm}
    \begin{nd}
        \have[n]{pq}{A \Conj B}
        \have[...]{}{\dots}
        \have[m]{}{B\qquad\qquad \ConjE,n}
    \end{nd}
\end{align*}
Namely, if at any prior point in the argument (say on line $n$) you have derived $A \Conj B$, then at any later point (say on line $m$) you can derive either $A$, or you can derive $B$.

Returning to the above argument, we use this rule again on line 4 (now justified by line 2). Having inferred both of the claims on lines 3 and 4, we next want to put them together to form our final conclusion. We do this with the \textbf{$\ConjTight$-Introduction Rule} justified by lines 3 and 4. Here is how it works in general:

\begin{center}
\begin{minipage}[t]{0.25\textwidth}
\textbf{$\ConjTight$-Introduction Rule:}
\vspace{-0.5cm}
\begin{align*}
     \begin{nd}
        \have[n]{}{A}
        \have[...]{}{\dots}
        \have[m]{}{B}
        \have[...]{}{\dots}
        \have[k]{}{A\Conj B\qquad\quad \ConjI,n,m}
    \end{nd}
\end{align*}
\end{minipage}
\hfill
\begin{minipage}[t]{0.19\textwidth}
Namely, if at some prior point in the argument (say on line $n$) you have derived $A$ and at some other point (say on line $m$) you have derived $B$, then at any later point (say on line $k$) you can derive $A \Conj B$.
\end{minipage}
\end{center}

By combining these two rules we have derived our desired conclusion $(q \Disj r) \Conj (s \Cond t)$ from our two premises. We denote this as:
$$p \Conj (q \Disj r),\quad (s \Cond t)\Conj p \quad \ProvesND \quad (q \Disj r) \Conj (s \Cond t)$$
Notice that someone could follow this argument perfectly well without knowing what the terms $\Disj$ and $\Cond$ mean.

An argument is \textbf{ND-valid} (denoted $X\ProvesND A$) if and only if the formula $A$ can be derived with $X$ as one's starting assumptions.

An argument is \textbf{ND-invalid} (denoted $X\not\ProvesND A$) if and only if there is no possible way to derive the formula $A$ with $X$ as one's starting assumptions. (In practice, there is no \textit{direct} way to show that an argument is ND-invalid: the space of all possible derivations is far too unwieldy to search by hand. There is, however, an \textit{indirect} way. Recall from above that natural deduction turns out to be sound (everything ND-derivable is table-valid), so if we can exhibit a truth-table countermodel witnessing $X \notEntails A$, it follows that $X \not\ProvesND A$ as well. We return to this in Lecture 12.)

\textbf{3. Natural Derivations using $\Disj$.} Here is a natural language example. Consider the following argument from premise $(p \Conj q) \Disj (p \Conj r)$ to conclusion $p$. Our premise tells us that either $p \Conj q$ or $p \Conj r$ holds, but we do not know which. Let's assume that we are in the first case, so that $p \Conj q$. It would then follow immediately that $p$. But now assume that we are in the second case, so that $p \Conj r$. It would then follow immediately that $p$. Hence whether we are in the first or the second case we have $p$. Therefore, $p$.

Here is how this argument would be written up as a \textbf{natural derivation}:
\begin{center}
\begin{minipage}[t]{0.2\textwidth}
As always, you begin by drawing a long vertical line (the main \textbf{scope line}). You then write a numbered list of formulas: these are your starting assumptions. You draw a horizontal stroke under your last assumption, to mark that everything above the stroke is \textit{assumed}, not derived. (Sometimes a derivation has no main assumptions at all.)
\end{minipage}
\hfill
\begin{minipage}[t]{0.1\textwidth}
\begin{align*}
\begin{nd}
    \hypo{1}{(p \Conj q) \Disj (p \Conj r)}
    \open
    \hypo{2}{p \Conj q}
    \have{3}{p \qquad \ConjE, 2}
    \close
    \open
    \hypo{4}{p \Conj r}
    \have{5}{p \qquad \ConjE, 4}
    \close
    \have{6}{p \qquad\quad \ \ \ \DisjE, 1, 2\text{--}3, 4\text{--}5}
\end{nd}
\end{align*}
\end{minipage}
\end{center}
Line 2 in this argument does not follow from any rule of inference. Instead, it is a \textbf{sub-assumption}, opening a \textbf{sub-derivation}. A sub-derivation gets its own, shorter scope line, nested inside the main one. Within this sub-assumption we can then apply any rules of inference that we like, making use of the new assumption. Here we derive line 3 from line 2 using the $\ConjE$ rule: namely, from $p\Conj q$ we infer $p$. Now that we have derived what we wanted from this sub-assumption we close it off. Beyond this point the sub-assumption is \textbf{discharged}: it is no longer available for use outside the sub-derivation. This is marked by the second scope line ending at line 3.

Line 4 of this argument then introduces another sub-assumption, namely $p \Conj r$. From here we then derive line 5 from line 4 using the $\ConjE$ rule: namely, from $p\Conj r$ we infer $p$. Having derived what we wanted from this sub-assumption, we close it off. This is marked by the second scope line ending at line 5.

In the final step of our argument we apply the \textbf{$\Disj$-Elimination Rule} justified by line 1 and the two sub-derivations 2-3 and 4-5. Here is how the rule works.

\textbf{$\Disj$-Elimination Rule (Proof by Cases):}
\vspace{-0.75cm}
\begin{align*}
\begin{nd}
\have[n]{}{A \Disj B}
\open
\hypo[n+1]{}{A}
\have[...]{}{\dots}
\have[m]{}{C}
\close
\open
\hypo[m+1]{}{B}
\have[...]{}{\dots}
\have[k]{}{C}
\close
\have[k+1]{}{C\qquad\qquad\DisjE, n, (n{+}1)\text{--}m, (m{+}1)\text{--}k}
\end{nd}
\end{align*}
\vspace{-0.75cm}\\
Namely, suppose that at some prior point in the argument (say on line $n$) you have derived $A\Disj B$. On line $n+1$ you then make a sub-assumption of $A$ and (somehow) argue your way to claim $C$ on line $m$. You then start another sub-argument leading from $B$ to $C$. Combining these two sub-arguments together gives the conclusion $C$ \textit{independent of our sub-assumptions about $A$ and $B$} (i.e., back on the main scope line).

Importantly, on our way from $B$ to $C$ we are allowed to invoke something on line $n$ or prior \textit{but not} our earlier assumption of $A$ (or anything between lines $n+1$ and $m$). But line $n$ is still fair game, since it sits outside that sub-derivation entirely. Here is the general rule for \textbf{accessibility}. A sub-derivation is \textit{closed off} once its scope line ends, and from that point on nothing written inside it may be cited again. So from a given line you may cite an earlier line only when every sub-derivation that the earlier line lives inside is still open at your current line, i.e., only when you can trace a path straight up the page to that earlier line without ever stepping across a scope line that has already ended. In particular, a formula on the main scope line is accessible everywhere below it, whereas a formula sitting inside a sub-derivation becomes inaccessible the instant that sub-derivation closes.

That was the elimination rule for $\Disj$, but what is its introduction rule? That is, what does it take to derive a claim of the form $A \Disj B$? It is deceptively simple. It follows immediately if we have derived $A$ earlier. Or if we derived $B$ earlier.

\textbf{$\Disj$-Introduction Rule:}
\vspace{-0.75cm}
\begin{align*}
    \begin{nd}
        \have[n]{}{A}
        \have[...]{}{\dots}
        \have[m]{}{A\Disj B\qquad\qquad \DisjI,n}
    \end{nd}
\hspace{2cm}
    \begin{nd}
        \have[n]{}{B}
        \have[...]{}{\dots}
        \have[m]{}{A\Disj B\qquad\qquad \DisjI,n}
    \end{nd}
\end{align*}
\vspace{-0.75cm}\\
Here is an example which demonstrates the power of these two introduction rules. Our previous argument established that $(p \Conj q) \Disj (p \Conj r) \ \ProvesND \ p$, but this conclusion seems weaker than it needs to be. The premise seems to tell us that (in addition to $p$ being true) either $q$ or $r$ is true. Let us therefore aim for a stronger conclusion: $(p \Conj q) \Disj (p \Conj r) \quad \ProvesND \quad p \Conj (q \Disj r)$.
\vspace{-0.75cm}
\begin{align*}
\begin{nd}
    \hypo{1}{(p \Conj q) \Disj (p \Conj r)}
    \open
    \hypo{2}{p \Conj q}
    \have{3}{p \qquad\qquad\qquad \ConjE, 2}
    \have{4}{q \qquad\qquad\qquad \ConjE, 2}
    \have{5}{q \Disj r \qquad\qquad \ \DisjI, 4}
    \have{6}{p \Conj (q \Disj r) \qquad \ConjI, 3, 5}
    \close
    \open
    \hypo{7}{p \Conj r}
    \have{8}{p \qquad\qquad\qquad \ConjE, 7}
    \have{9}{r \qquad\qquad\qquad \ConjE, 7}
    \have{10}{q \Disj r \qquad\qquad \ \DisjI, 9}
    \have{11}{p \Conj (q \Disj r) \qquad \ConjI, 8, 10}
    \close
    \have{12}{p \Conj (q \Disj r) \qquad\qquad \DisjE, 1, 2\text{--}6, 7\text{--}11}
\end{nd}
\end{align*}
As before, this is a proof by cases. Lines 5 and 10 are strict weakenings of lines 4 and 9 respectively: we go from $q$ to $q\Disj r$ and from $r$ to $q\Disj r$. The reason for retreating to these weakened claims in each case is that \textit{this weakened claim is precisely what the two cases agree about}. Indeed, lines 6 and 11 reach the same conclusion $p \Conj (q \Disj r)$ and hence (by cases) the conclusion follows.

We have improved our argument by strengthening its conclusion from $p$ to the stronger $p \Conj (q \Disj r)$:
$$(p \Conj q) \Disj (p \Conj r) \quad \ProvesND \quad p \Conj (q \Disj r)$$
In fact, as we shall now see, this last derivation is \textit{lossless}: the conclusion is just as strong as the premise. Concretely, we can reverse it as follows:
$$p \Conj (q \Disj r) \quad \ProvesND \quad (p \Conj q) \Disj (p \Conj r)$$
That is, beginning from $p \Conj (q \Disj r)$ we can apply our natural deduction rules to reach $(p \Conj q) \Disj (p \Conj r)$. Here is the proof.
\vspace{-0.75cm}
\begin{align*}
\begin{nd}
    \hypo{1}{p \Conj (q \Disj r)}
    \have{2}{p \qquad\qquad\qquad\qquad\qquad \ \ConjE, 1}
    \have{3}{q \Disj r \qquad\qquad\qquad\qquad\quad \ConjE, 1}
    \open
    \hypo{4}{q}
    \have{5}{p \Conj q \qquad\qquad\qquad\quad \ConjI, 2, 4}
    \have{6}{(p \Conj q) \Disj (p \Conj r) \qquad \DisjI, 5}
    \close
    \open
    \hypo{7}{r}
    \have{8}{p \Conj r \qquad\qquad\qquad\quad \ \ConjI, 2, 7}
    \have{9}{(p \Conj q) \Disj (p \Conj r) \qquad \DisjI, 8}
    \close
    \have{10}{(p \Conj q) \Disj (p \Conj r) \qquad\qquad \DisjE, 3, 4\text{--}6, 7\text{--}9}
\end{nd}
\end{align*}
Thus, the claims $p \Conj (q \Disj r)$ and $(p \Conj q) \Disj (p \Conj r)$ are \textit{inter-derivable}: We can get from one to the other and vice versa using natural deduction rules. This gives us a notion of logical equivalence for natural deduction.

Formulas $A$ and $B$ are \textbf{ND-equivalent} if and only if $A\ProvesND B$ and $B\ProvesND A$.

Contrast this with the notion of equivalence coming from our discussion of truth tables in Lecture 4: $A$ and $B$ are \textit{table-equivalent} if and only if they take on the same truth-values in all models (i.e., at every row in their truth tables). Also contrast this with the notion of equivalence coming from our discussion of trees in Lecture 6: $A$ and $B$ are \textit{tree-equivalent} if and only if $A\Bicond B$ is a tree-tautology.

To deepen our comparison with the tree method, we shall next discuss how natural deduction handles the biconditional, $\Bicond$, and tautologies (and contradictions and contingent statements).

\textbf{4. Natural Derivations using $\Bicond$.} Next up is the biconditional. What are its rules of inference? If we know $A\Bicond B$ and we know $A$ then we can conclude $B$. Similarly, if we know $A\Bicond B$ and we know $B$ then we can conclude $A$.

\textbf{$\Bicond$-Elimination Rule:}
\vspace{-0.5cm}
\begin{align*}
 \begin{nd}
        \have[n]{}{A \Bicond B}
        \have[...]{}{\dots}
        \have[m]{}{A}
        \have[...]{}{\dots}
        \have[k]{}{B\qquad\qquad\BicondE,n,m}
    \end{nd}
    \hspace{1cm}
 \begin{nd}
        \have[n]{}{A \Bicond B}
        \have[...]{}{\dots}
        \have[m]{}{B}
        \have[...]{}{\dots}
        \have[k]{}{A\qquad\qquad\BicondE,n,m}
    \end{nd}
\end{align*}
But how could we establish that $A\Bicond B$ in the first place? To do this one just has to show that we can derive $B$ from $A$ and that we can derive $A$ from $B$.

\textbf{$\Bicond$-Introduction Rule:}
\vspace{-1cm}
\begin{align*}
\begin{nd}
\open
\hypo[n]{1}{A}
\have[...]{}{\dots}
\have[m]{2}{B}
\close
\open
\hypo[m+1]{3}{B}
\have[...]{}{\dots}
\have[k]{4}{A}
\close
\have[k+1]{5}{A \Bicond B\qquad\qquad \BicondI, n\text{--}m, (m{+}1)\text{--}k}
    \end{nd}
\end{align*}
\vspace{-0.75cm}\\
Let's see this introduction rule in action. We can combine our proof of $(p \Conj q) \Disj (p \Conj r) \ \ProvesND \ p \Conj (q \Disj r)$ with our proof of $p \Conj (q \Disj r) \ \ProvesND \ (p \Conj q) \Disj (p \Conj r)$ as follows:
\vspace{-0.75cm}
\begin{align*}
\begin{nd}
    \open
    \hypo{1}{(p \Conj q) \Disj (p \Conj r)}
    \have[\vdots]{}{\vdots \qquad \text{(Insert the first proof here)}}
    \have[12]{12}{p \Conj (q \Disj r)}
    \close
    \open
    \hypo[13]{13}{p \Conj (q \Disj r)}
    \have[\vdots]{}{\vdots \qquad \text{(Insert the second proof here)}}
    \have[22]{22}{(p \Conj q) \Disj (p \Conj r)}
    \close
    \have[23]{23}{((p \Conj q) \Disj (p \Conj r)) \Bicond (p \Conj (q \Disj r)) \qquad \BicondI, 1\text{--}12, 13\text{--}22}
\end{nd}
\end{align*}
\vspace{-0.75cm}\\
Interestingly this argument has no main assumptions! The first line immediately introduces a sub-assumption rather than introducing the argument's premises. This is not a typo: \textit{this argument has no premises}. We denote this as: $\ProvesND \quad ((p \Conj q) \Disj (p \Conj r)) \Bicond (p \Conj (q \Disj r))$.

A formula $A$ is an \textbf{ND-tautology} (denoted $\ProvesND A$) if and only if it can be derived from no assumptions. (Compare this with our table and tree notions.)

\textbf{5. Next Time...} There are still a few more natural deduction rules to introduce (namely, for $\Neg$ and $\Cond$). These are (somewhat) controversial and I am prepared to teach the controversy!
\end{document}
